\documentclass{article}

% packages
\usepackage{amsmath, amsthm, thmtools, amsfonts, amssymb, luacode, catchfile, tikzducks, hyperref, ifthen}
\ifcsname c@kobocompile\endcsname
	\usepackage[a5paper, total={1072pt, 1448pt}, margin=10pt, includeheadfoot]{geometry} % set page margins
\else
	\usepackage[a4paper, margin=50pt, includeheadfoot]{geometry}
\fi
\usepackage[shortlabels]{enumitem}
\usepackage[skip=3pt, indent=0pt]{parskip}

% language
\usepackage[bidi=basic, layout=tabular, provide=*]{babel}
\ifcsname c@english\endcsname
	\babelprovide[main, import]{english}
\else
	\babelprovide[main, import]{hebrew}
	\babelprovide{rl}
\fi
%\babelfont{rm}{Libertinus Serif}
\babelfont{rm}[Renderer=Harfbuzz]{Libertinus Serif}
\babelfont{sf}{Libertinus Sans}
\babelfont{tt}{Libertinus Mono}

% style
\AddToHook{cmd/section/before}{\clearpage}	% Add line break before section
\linespread{1.3}
\setcounter{secnumdepth}{0}		% Remove default number tags from sections, this won't do well with theorems
\AtBeginDocument{\setlength{\belowdisplayskip}{3pt}}
\AtBeginDocument{\setlength{\abovedisplayskip}{3pt}}
\graphicspath{ {../images/} }

% operators
\DeclareMathOperator\cis{cis}
\DeclareMathOperator\Sp{Sp}
\DeclareMathOperator\tr{tr}
\DeclareMathOperator\im{Im}
\DeclareMathOperator\re{Re}
\DeclareMathOperator\diag{diag}
\DeclareMathOperator*\lowlim{\underline{lim}}
\DeclareMathOperator*\uplim{\overline{lim}}
\DeclareMathOperator\rng{rng}
\DeclareMathOperator\Sym{Sym}
\DeclareMathOperator\Arg{Arg}
\DeclareMathOperator\Log{Log}
\DeclareMathOperator\dom{dom}
\DeclareMathOperator\supp{Supp}
\DeclareMathOperator\var{Var}
\DeclareMathOperator\cov{Cov}

% commands
%\renewcommand\qedsymbol{\textbf{מש''ל}}
%\renewcommand\qedsymbol{\fbox{\emoji{lizard}}}
\newcommand{\Aa}[0]{\mathcal{A}}
\newcommand{\Bb}[0]{\mathcal{B}}
\newcommand{\CC}[0]{\mathbb{C}}
\newcommand{\Cc}[0]{\mathcal{C}}
\newcommand{\EE}[0]{\mathbb{E}}
\newcommand{\FF}[0]{\mathbb{F}}
\newcommand{\Ff}[0]{\mathcal{F}}
\newcommand{\Ii}[0]{\mathcal{I}}
\newcommand{\Gg}[0]{\mathcal{G}}
\newcommand{\Ll}[0]{\mathcal{L}}
\newcommand{\Mm}[0]{\mathcal{M}}
\newcommand{\NN}[0]{\mathbb{N}}
\newcommand{\Nn}[0]{\mathcal{N}}
\newcommand{\PP}[0]{\mathbb{P}}
\newcommand{\Pp}[0]{\mathcal{P}}
\newcommand{\QQ}[0]{\mathbb{Q}}
\newcommand{\RR}[0]{\mathbb{R}}
\newcommand{\Rr}[0]{\mathcal{R}}
\newcommand{\Ss}[0]{\mathcal{S}}
\newcommand{\TT}[0]{\mathbb{T}}
\newcommand{\Uu}[0]{\mathcal{U}}
\newcommand{\Vv}[0]{\mathcal{V}}
\newcommand{\Ww}[0]{\mathcal{W}}
\newcommand{\ZZ}[0]{\mathbb{Z}}
\newcommand{\acts}[0]{\circlearrowright}
\newcommand{\explain}[2] {
	\begin{flalign*}
		 && \text{#2} && \text{#1}
	\end{flalign*}
}
\newcommand{\maketitleprint}[0]{ \begin{center}
	%\begin{tikzpicture}[scale=3]
	%	\duck[graduate=gray!20!black, tassel=red!70!black]
	%\end{tikzpicture}	
	\includegraphics[width=6cm]{cover}
\end{center}
}

% theorem commands
\newtheoremstyle{c_remark}
	{}	% Space above
	{}	% Space below
	{}% Body font
	{}	% Indent amount
	{\bfseries}	% Theorem head font
	{}	% Punctuation after theorem head
	{.5em}	% Space after theorem head
	{\thmname{#1}\thmnumber{ #2}\thmnote{ \normalfont{\text{(#3)}}}}	% head content
\newtheoremstyle{c_definition}
	{3pt}	% Space above
	{3pt}	% Space below
	{}% Body font
	{}	% Indent amount
	{\bfseries}	% Theorem head font
	{}	% Punctuation after theorem head
	{.5em}	% Space after theorem head
	{\thmname{#1}\thmnumber{ #2}\thmnote{ \normalfont{\text{(#3)}}}}	% head content
\newtheoremstyle{c_plain}
	{3pt}	% Space above
	{3pt}	% Space below
	{\itshape}% Body font
	{}	% Indent amount
	{\bfseries}	% Theorem head font
	{}	% Punctuation after theorem head
	{.5em}	% Space after theorem head
	{\thmname{#1}\thmnumber{ #2}\thmnote{ \text{(#3)}}}	% head content

\ifcsname c@english\endcsname
	\theoremstyle{plain}
	\newtheorem{theorem}{Theorem}[section]
	\newtheorem{lemma}[theorem]{Lemma}
	\newtheorem{proposition}[theorem]{Proposition}
	\newtheorem*{proposition*}{Proposition}
	%\newtheorem{corollary}[theorem]{אין חלופה עברית}

	\theoremstyle{definition}
	\newtheorem{definition}[theorem]{Definition}
	\newtheorem*{definition*}{Definition}
	\newtheorem{example}{Example}[section]
	\newtheorem{exercise}{Exercise}[section]

	\theoremstyle{remark}
	\newtheorem*{remark}{Remark}
	\newtheorem*{solution}{Solution}
	\newtheorem{conclusion}[theorem]{Conclusion}
	\newtheorem{notation}[theorem]{Notation}
\else
	\theoremstyle{c_plain}
	\newtheorem{theorem}{משפט}[section]
	\newtheorem{lemma}[theorem]{למה}
	\newtheorem{proposition}[theorem]{טענה}
	\newtheorem*{proposition*}{טענה}
	%\newtheorem{corollary}[theorem]{אין חלופה עברית}

	\theoremstyle{c_definition}
	\newtheorem{definition}[theorem]{הגדרה}
	\newtheorem*{definition*}{הגדרה}
	\newtheorem{example}{דוגמה}[section]
	\newtheorem{exercise}{תרגיל}[section]

	\theoremstyle{c_remark}
	\newtheorem*{remark}{הערה}
	\newtheorem*{solution}{פתרון}
	\newtheorem{conclusion}[theorem]{מסקנה}
	\newtheorem{notation}[theorem]{סימון}
\fi

% Questions related commands
\newcounter{question}
\setcounter{question}{1}
\newcounter{sub_question}
\setcounter{sub_question}{1}

\ifcsname c@english\endcsname
	\newcommand{\question}[1][0]{
		\ifthenelse{#1 = 0}{}{\setcounter{question}{#1}}
		\section{Question \arabic{question}}
		\addtocounter{question}{1}
		\setcounter{sub_question}{1}
	}

	\newcommand{\subquestion}[1][0]{
		\ifthenelse{#1 = 0}{}{\setcounter{sub_question}{#1}}
		\subsection{Part \alph{sub_question}}
		\addtocounter{sub_question}{1}
	}
\else
	\newcommand{\question}[1][0]{
		\ifthenelse{#1 = 0}{}{\setcounter{question}{#1}}
		\section{שאלה \arabic{question}}
		\addtocounter{question}{1}
		\setcounter{sub_question}{1}
	}

	\newcommand{\subquestion}[1][0]{
		\ifthenelse{#1 = 0}{}{\setcounter{sub_question}{#1}}
		\subsection{סעיף \localecounter{letters.gershayim}{sub_question}}
		\addtocounter{sub_question}{1}
	}
\fi

% import lua and start of document
\directlua{common = require ('../common')}

\GetEnv{AUTHOR}

% headers
\author{\AUTHOR}
\date\today

\title{פתרון מטלה 07 --- אנליזה על יריעות, 80426}

\DeclareMathOperator{\vol}{vol}

\begin{document}
\maketitle
\maketitleprint[blue]

\question{}
ניזכר כי $TM$ הוא האגד המשיק של היריעה $M$.
תהי $\pi : TM \to M$ ההעתקה $(x, v) \mapsto x$, ונגדיר $s : M \to TM$ על־ידי $s(x) = (x, 0)$.
תהי $f : M \to N$ חלקה, אז נקבל ש־$df : TM \to TN$ העתקה המוגדרת על־ידי $df(x, v) = (f(x), Df |_x(v))$.

\subquestion{}
נראה ש־$\pi, s$ חלקות.
\begin{proof}
	נבחין כי שתי הפונקציות הן לינאריות ולכן בפרט חלקות, אבל נוכיח את הטענה ישירות מהגדרה.

	נניח ש־$M^k \subseteq \RR^m$.
	אז $TM \subseteq \RR^{m + k}$.
	תהי $(p, v) \in TM$ ונבחן את $\overline{\pi} : \RR^{m + k} \to \RR^m$ הרחבה רציפה של $\pi : TM \to M$.
	כלומר, נגדיר ש־$\overline{\pi} \restriction TM = \pi$, וכן $\overline{\pi}(x, u) = x$, זוהי העתקה לינארית $\RR^{m + k} \to \RR^m$, כלומר העתקה לינארית אוקלידית, ולכן היא חלקה.

	נבצע תהליך דומה ל־$s$, נגדיר $\overline{s} : \RR^m \to \RR^{m + k}$ על־ידי $\overline{s}(x) = (x, 0)$, אז בבירור $\overline{s} \restriction M = s$, ו־$\overline{s}$ חלקה ומעידה על $s$ כהעתקה חלקה.
\end{proof}

\subquestion{}
נראה ש־$df$ היא חלקה.
\begin{proof}
	נניח ש־$N^l \subseteq \RR^n$. אז $f \subseteq {(\RR^n)}^{\RR^m}$ וכן $D f |_x : T_x(M) \to T_{f(x)}(N)$, כלומר $Df |_x \subseteq {(\RR^l)}^{\RR^k}$, לכל $x \in M$. \\
	הפעם נרצה לבנות את ההרחבה כך שקיימת פונקציה חלקה $\overline{df} : U \to \RR^{n + l}$ כך ש־$U \subseteq \RR^{m + k}$ פתוחה, וכן ש־$\overline{df} \restriction TM = df$.
	נניח ש־$(p, v) \in TM$ ונרצה להראות ש־$df$ חלקה ב־$(x, v)$.
	ידוע כי $f$ חלקה ולכן קיימת $\overline{f} : U_0 \to \RR^n$ חלקה, עבור $p \in U_0 \subseteq \RR^m$ פתוחה.
	לכל $x \in U_0$ נבחין כי $D \overline{f} |_x$ היא העתקה לינארית ולכן חלקה אף היא.
	נגדיר אם כך את $\overline{df}(x, v) = (\overline{f}(x), D \overline{f}|_x(v))$ ונקבל שזו העתקה חלקה, אבל $\overline{df} \restriction TM = df$ ישירות מהגדרת הנגזרת על יריעה, ולכן פונקציה זו מעידה כי $df$ בעצמה חלקה.
\end{proof}

\question{}
ניזכר כי,
\[
	\operatorname{SL}_n(\RR)
	= \{ A \in M_n(\RR) \mid \det(A) = 1 \}
\]
היא קבוצת ההעתקות הלינאריות $\RR^n \to \RR^n$ ההפיכות והאורתוגונליות.

\subquestion{}
נראה ש־$\operatorname{SL}_n(\RR)$ היא יריעה.
\begin{proof}
	נגדיר $f : \operatorname{GL}_n(\RR) \to \RR$ על־ידי $f(A) = \det(A)$.
	\begin{align*}
		D f |_{I_n}(B)
		& = \lim_{h \to 0} \frac{f(I_n + h B) - f(I_n)}{h} \\
		& = \lim_{h \to 0} \frac{\det(I_n + h B) - \det(I_n)}{h} \\
		& = \lim_{h \to 0} \frac{(1 + h B_{1, 1}) \det(I_{n - 1} + h B_{i > 1, j > 1}) + h(\ldots) - 1}{h} \\
		& = \lim_{h \to 0} \frac{(1 + h B_{1, 1}) \cdots (1 + h B_{n, n}) - 1}{h} + o(h) \\
		& = \lim_{h \to 0} \frac{1 + h \tr(B) - 1}{h} + o(h) \\
		& = \tr(B)
	\end{align*}
	נבחין כי $D f |_{I_n}$ היא על, נוכל לבחור $\diag(r, 1, \ldots, 1)$ לכל $r \in \RR$, ולכן $I_n$ נקודה רגולרית.
	זה לא מוכיח כלום.
\end{proof}

\subquestion{}
נחשב את $\dim \operatorname{SL}_n(\RR)$.
\begin{solution}
	TODO
\end{solution}

\question[4]
נגדיר $M_a^2 = \{ (x, y, z, w) \in \RR^4 \mid x^2 + y^2 = 1, z^2 + w^2 = {(x + a)}^2 \}$.

\subquestion{}
נראה ש־$M_3^2$ היא יריעה.
\begin{proof}
	נגדיר את הפונקציה $g : \RR^4 \to \RR^2$ על־ידי,
	\[
		g(x, y, z, w)
		= (x^2 + y^2, z^2 + w^2 - {(x + a)}^2)
	\]
	אז נקבל ש־$g^{-1}(\{ (1, 0) \}) = M_a^2$.
	ממשפט הפונקציה הסתומה ליריעות מספיק שנוכיח ש־$(1, 0)$ ערך רגולרי של $f$.
	נגזור את הפונקציה,
	\[
		D g |_{(x, y, z, w)}
		= \begin{pmatrix}
			2x & 2y & 0 & 0 \\
			-2(x + a) & 0 & 2z & 2w
		\end{pmatrix}
	\]
	ונקבל כי לכל הצבת ערך יש לפחות שתי שורות בלתי תלויות־לינארית, כלומר הנקודה רגולרית לכל נקודה, \\
	ובפרט אם $(x, y, z, w) \in f^{-1}(\{ (1, 0) \})$.
	נסיק ממשפט הפונקציה הסתומה שאכן $M_3^2$ יריעה.
\end{proof}

\subquestion{}
נראה ש־$M_0^2$ איננה יריעה.
\begin{proof}
	נניח בשלילה ש־$M_0^2$ יריעה, ונגדיר,
	\[
		f : M \to S^1,
		\qquad
		f(x, y, z, w) = (x, y)
	\]
	זוהי העתקה חלקה כהעתקת צמצום ולכל $q \in S^1$ ערך רגולרי גם $f^{-1}(q)$ יריעה.
	אבל מהגדרת $f$ נקבל שכל ערך ב־$S^1$ היא רגולרית, זאת שכן,
	\[
		D f \mid_{(x, y, z, w)}
		= \begin{pmatrix}
			1 & 0 & 0 & 0 \\
			0 & 1 & 0 & 0
		\end{pmatrix} 
	\]
	ולכן בפרט $f^{-1}(\{(0, 1)\})$ יריעה, ממימד 2.
	אבל,
	\[
		f(x, y, z, w) = (0, 1)
		\iff x^2 + y^2 = 1, w^2 + z^2 = x^2, x = 0, y = 1
		\iff x = 0, y = 1, z = 0, w = 0
	\]
	כלומר זוהי יריעה ממימד 0, בסתירה למשפט הפונקציה הסתומה.
\end{proof}

\question{}
נגדיר את העתקת הופ כ־$h : S^3(1) \to S^2(\frac{1}{2})$ המוגדרת על־ידי,
\[
	h(z, w)
	= (z \overline{w}, \frac{1}{2} ({|z|}^2 - {|w|}^2))
\]
נראה כי אכן תמונתה היא ב־$S^2(\frac{1}{2})$ ושכל $q \in S^2(\frac{1}{2})$ היא ערך רגולרי של $h$.
\begin{proof}
	נבחין כי $\frac{1}{2} ({|z|}^2 - {|w|}^2) \in \RR$ וכן ש־$z \overline{w} \in \CC$ ולכן $h$ משוכנת ב־$\RR^3$.
	עוד נבחין כי,
	\[
		{|h(z, w)|}^2
		= z \overline{w} w \overline{z} + \frac{1}{4} {|z|}^4 - \frac{1}{2} {|zw|}^2 + \frac{1}{4} {|w|}^4
		= \frac{1}{4} {|z|}^4 + \frac{1}{2} {|zw|}^2 + \frac{1}{4} {|w|}^4
		= {\left(\frac{1}{2} ({|z|}^2 + {|w|}^2)\right)}^2
		= {\left(\frac{1}{2} \cdot 1\right)}^2
	\]
	ונובע שתמונת $h$ היא ב־$S^2(\frac{1}{2})$ בלבד.

	נעבור לבדיקת רגולריות.
	נגזור את איברי המספרים המרוכבים,
	\[
		D h |_{(z, w)}
		= D h |_{(z_r, z_i, w_r, w_i)}
		= \begin{pmatrix}
			w_r & - w_i & z_r & -z_i \\
			w_i & w_r & z_i & z_r \\
			\frac{1}{2} z_r & \frac{1}{2} z_i & \frac{1}{2} w_r & \frac{1}{2} w_i
		\end{pmatrix} 
		= \begin{pmatrix}
			\overline{w} & \overline{z} \\
			i \overline{w} & i \overline{z} \\
			\frac{1}{2} z & \frac{1}{2} w
		\end{pmatrix} 
	\]
	והנקודה $(z, w)$ רגולרית אם ההעתקה היא על $\RR^2 \simeq \CC \simeq T_{h(z, w)}(S^2(\frac{1}{2}))$.
	אבל נבחין כי עמודותיה בלתי־תלויות לינארית עבור כל $(z, w) \ne 0$, וכן $0 \notin S^3(1)$, ולכן נסיק שהיא רגולרית בכל נקודה.
	נובע אם כך שלכל $q \in S^2(\frac{1}{2})$ אכן כל מקורותיה רגולריים וכן היא ערך רגולרי.
\end{proof}

\end{document}
